\chapter{Results}
\chaptermark{} %this is the chapter heading that will show on subsequent pages

\section{Computational Model and DA Framework}

\subsection{Core framework}



The first output of this work is a general data assimilation framework for MATLAB.
By utilizing an object-oriented (OO) design, the model and data assimilation algorithm code are separate and can be changed independently.
The principal advantage of this approach is the ease of incorporation of new models and DA techniques.

To add a new model to this framework, it is only necessary to write a MATLAB class wrapper for the model, and the existing Lorenz 63 and OpenFOAM models are templates.
Testing a new technique for performing the DA step is as simple as writing a MATLAB function for the assimilation steps, with the model forecast and observations as input.

The core of this framework is the MATLAB script modelDAinterface.m which wraps observations, the model class, and any DA method together.
A description of the main algorithm is provided in Appendix \ref{app:foamlab}.

For large models, it is necessary to perform the assimilation in a local space, and this consideration is accounted for within the framework.
This procedure, referred to as localization, eliminates spurious long-distance correlations and parallelizes the assimilation.

\begin{figure}[t!]
  \centering
  \includegraphics[width=0.79\textwidth]{/Users/andyreagan/work/2013/2013-05data-assimilation/presentation/foamLab_slides001.pdf}
  \caption{
    modelDAinterface.m schematic
  }
  \label{fig:modelDAschem}
\end{figure}

\subsection{Tuned OpenFOAM model}

In addition to wrapping the OpenFOAM model into MATLAB, defining a physically realistic simulation was necessary for accurate results.
A non-uniform sampling of the OpenFOAM model parameter space across BC type, BC values, and turbulence model was performed.
The amount of output data constitutes what is today ``big data,'' and for this reason it was difficult to find a suitable parameter regime systematically.
Specifically, each model run generates 5MB of data per time save, and with 10000 timesteps this amounts to roughly 50GB of data for each model run.
We reduce the data to time-series for a subset of model variables, namely timeseries of the reconstructed mass flow rate and temperature difference $\Delta T_{3-9}$.
From these timeseries, we then select the BC, BC values, and turbulence model empirically based on the flow behavior observed.

\begin{figure}[h!]
  \centering
  \includegraphics[width=0.49\textwidth]{/Users/andyreagan/work/2013/2013-05data-assimilation/presentation/blank.pdf}
  \caption{
    mesh verification
  }
  \label{fig:meshverification}
\end{figure}

\section{Prediction Skill for Synthetic Data}

We first present the results pertaining to the accuracy of forecasts for synthetic data.
There are many possible experiments given the choice of assimilation window, data assimilation algorithm, localization scheme, model resolution, observational density, observed variables, and observation quality.
We focus on considering the effect of observations and observational locations on the resulting forecast skill.

\subsubsection{Covariance Localization}

\begin{figure}[h!]
  \centering
  \includegraphics[width=0.79\textwidth]{/Users/andyreagan/work/2013/2013-05data-assimilation/presentation/foamLab_slides002.pdf}
  \caption{
    covariance localization zones
  }
  \label{fig:localcovzones}
\end{figure}

\begin{figure}[h!]
  \centering
  \includegraphics[width=0.79\textwidth]{/Users/andyreagan/work/2013/2013-05data-assimilation/presentation/foamLab_slides003.pdf}
  \caption{
    covariance localization radii
  }
  \label{fig:localcovradii}
\end{figure}

The need for covariance localization was discussed in Chapter 2, and here we implement two different schemes.
First, we consider cells within each perpendicular slice of the loop to be a cohesive group, and their neighbors those cells within $r$ slices (where here, $r$ is an integer).
Computationally this scheme is the simplest to implement, and results in the fewest covariance localizations.

Next, we consider localization by distance from center.
Data assimilation is performed for each grid point, considering neighboring cells and observations with cell center within a radius $r$ from each point (e.g. $r = 0.5$cm).

\subsubsection{Generation of Synthetic Data}

Synthetic data is generated from one model run, with a modified ``buoyantBoussinesqPimpleFoam'' solver applied to fixed value boundary conditions of 290K top and 340K bottom, 80,000 grid points, a time step of .01, and laminar turbulence model.
Data is saved every half second, and is subsampled for 2,4,8,16 and 32 synthetic temperature sensors spaced evenly around the loop.
The temperature at these synthetic sensors is reported as the average for cells with centers within 0.5cm of the observation center.
Going beyond 32 temperature observations, which we find to be necessary, samples are performed by observed cell-center temperature values.
The position of the chosen cell depends on the number of observations being made, and within that number there are options that are tested.
For less than 1000 observations, (1000 being the number of angular slices) observations are placed at slice centers, for slices spaced maximally apart.
For 1000 observations both slice center (midpoint going all the way around) and staggered (closer to one wall, alternating) are implemented.
%% For greater than 1000 observations, we space take observations at locations maximally spaced within each slice, for each slice.
For $n$ observations with $n>1000$, assuming that $n/1000$ is an integer, for each slice we space the $n/1000$ observations maximally far apart within each slice.

In addition to temperature sensing, the velocity in $y$ and $z$ is sampled at 50, 100, and 300 points chosen randomly at each time step.
These velocity observations are designed to simulate a reconstructed velocity measurement based on video particle tracking.
The following results do not include the use of these velocity observations.
To obtain realistic sampling intervals, this data was aggregated and is then subsampled at the assimilation window length.

For both the Lorenz 63 model and OpenFOAM, error from the observed $T_{3-9}$ is computed as RMSE, such that the models can be compared.
In the EM model, a maximum of 2 temperature sensors are considered for observations at $T_3$ and $T_9$, and no mean velocity reconstruction from subsampled velocity is attempted.

For any of the predictions to work (even with complete observation coverage), it was first necessary to scale the model variables to the same order of magnitude.
This is accomplished here by dividing through by the mean climatological values from the observations and model results before assimilation (e.g. multiplying temperature by $1/300$). 
While other, more complex, strategies may result in better performance there are no general rules to determine the best scaling factors and good scaling is often problem dependent \shortcite{navon1992variational,Mitchell2012PhD}.

\subsection{Observational Network}

In general, we see that increasing observational density leads to improved forecast accuracy.
With too few observations, the data assimilation is unable to recover the underlying dynamics.

\begin{figure}[h!]
  \centering
  \includegraphics[width=0.79\textwidth]{/Users/andyreagan/work/2013/2013-05data-assimilation/src/experiments/openFoamAssimilation/data/synData-tStep.01-b290-t340-m40000/results_vs_N_w10.png}
  \caption{
    Prediction skill over time for synthetic temperature sensors.
    With these observations, our DA scheme is unable to predict the model run at all.
    Compared with operational NWP, we are observing far less model variables in this situation, with a $10^5$ times more model variables than observations, even for the 32 sensors.
  }
  \label{fig:TsensorExp}
\end{figure}

While the inability to predict the flow direction with these synthetic temperature sensors may cast doubt on the possibility of predicting a physical experiment using CFD, we discuss improvements to our prediction system that may be able to overcome these initial difficulties.
Regardless, in Figure \ref{fig:fullObsExp} we verify directly that our assimilation procedure is working adequately with a sufficient number of observations: temperature at all cells.

\begin{figure}[h!]
  \centering
  \includegraphics[width=0.79\textwidth]{/Users/andyreagan/work/2013/2013-05data-assimilation/src/experiments/openFoamAssimilation/data/synData-tStep.01-b290-t340-m40000/results_vs_N_w10_40000.png}
  \caption{
    Prediction skill over time for complete observations of temperature.
    With these observations, our DA scheme is able to predict the model run well.
    Increasing the number of ensemble members $N$ leads to a decrease in the forecast error, as one would hope with a more complete sampling of model uncertainty.
    Note that although we may full observations of temperature, it is only one of five model variables assimilated.
  }
  \label{fig:fullObsExp}
\end{figure}

With success at predicting flow direction given full temperature observations, we now ask the question: just how many observations are necessary?
The aforementioned generation of observational data considers the different possible observational network approaches for given number of observations $N_\text{obs}$.
From these approaches, we select the best method available to assimilation $N_\text{obs}$ for each $N$.
Our results initially indicate that the averaging of temperature over many cells for the synthetic temperature sensors up to 32 sensors degraded prediction skill.
\todo{still pulling all these results together}

\begin{figure}[h!]
  \centering
  \includegraphics[width=0.49\textwidth]{/Users/andyreagan/work/2013/2013-05data-assimilation/presentation/blank.pdf}
  \caption{
    skill vs $N_\text{obs}$
  }
  \label{fig:obsexp}
\end{figure}



%% And finally, we attempt a novel adaptive covariance localization scheme based on the identifying coherent flow structures.
%% Flow structures are identified by considering the leading vectors in $U$ of the singular value decomposition (SVD) of flow within a reasonably sized domain (e.g. within slice localization regions for large $r$).
%% This corresponds to a proper orthogonal decomposition (POD) of the flow structures.

\section{Conclusion}

Throughout the process of building a data assimilation scheme from scratch, as well as deriving and tuning a computational model, we have learned a great deal about the difficulties and successes of modeling.
We have found that data assimilation algorithms and CFD solvers are sensitive to parameter choices, and cannot be used successfully without an approach that is specific to the problem at hand.
Such an approach has been developed, and proved to be successful.

For the future of data assimilation, the numerical schemes will need to capitalize on new computing resources.
The trend toward massively parallel computation puts the spotlight on effective localization of the assimilation step, which may be difficult in schemes that do not use ensembles, such as the 4D-Var.
Maintaining TLM and adjoint code is difficult and time-consuming, and with future increases in model resolution and atmospheric understanding, TLMs of full GCM models may become impossible.
Again, this points to the potential use of local ensemble based filters such as the LETKF as implemented here.

\subsection{Future directions}

The immediate next step for this work is to test prediction skill in real time.
All of the pieces are there: an observation-based data assimilation framework and models that run at near real-time speed.
Currently on the VACC, the OpenFOAM simulation with grid size 36K runs with approximately half of real-time speed on one core.
This allows the possibility of running $N$ ensemble methods using $4N$ cores, allowing time for the assimilation.

\begin{figure}[h!]
  \centering
  \includegraphics[width=0.79\textwidth]{/Users/andyreagan/work/2013/2013-05data-assimilation/presentation/foamLab_slides004.pdf}
  \caption{
    covariance localization adaptive
  }
  \label{fig:localcovadaptive}
\end{figure}

Beyond real-time prediction, it is now possible to design and test new methods of data assimilation for large scale systems.
New methods specific to CFD, possibly using principle mode decomposition of flow fields have the potential to greatly improve the skill of real-time CFD applications.
In Figure \ref{fig:localcovadaptive}, flow structures present in our computational experiment demonstrate the potential for adaptive covariance localization.

The numerical coupling of CFD to experiment by DA should be generally useful to improve the skill of CFD predictions in even data-poor experiments, which can provide better knowledge of unobservable quantities of interest in fluid flow.

%% \begin{center}
%% In data we trust. 
%% \end{center}


